\documentclass[11pt,a4paper]{article}
\usepackage[margin=21mm]{geometry}
\usepackage[T1]{fontenc}\usepackage{lmodern}
\usepackage{amsmath,amssymb,booktabs,array,tikz,hyperref}
\usetikzlibrary{calc,arrows.meta}
\setlength{\parindent}{0pt}\setlength{\parskip}{5pt}
\newcommand{\VDD}{V_{\mathrm{DD}}}
% Vertical insulated-gate switch symbol; pMOS has a gate bubble.
% The final argument selects a bold channel for high-Vt devices.
\newcommand{\MOS}[6]{
 \coordinate (#1t) at (#2,#3+.45);\coordinate (#1b) at (#2,#3-.45);
 \draw (#1t)--(#2,#3+.25)--(#2-.13,#3+.25);
 \draw[line width=#6] (#2-.13,#3+.27)--(#2-.13,#3-.27);
 \draw (#2-.13,#3-.25)--(#2,#3-.25)--(#1b);
 \draw (#2-.29,#3+.25)--(#2-.29,#3-.25);
 \def\kind{#5}\def\pkind{p}
 \ifx\kind\pkind \draw (#2-.39,#3) circle(.05);\draw (#2-.44,#3)--(#2-.75,#3);
 \else \draw (#2-.29,#3)--(#2-.75,#3);\fi
 \node[left,font=\scriptsize] at (#2-.75,#3) {$#4$};
}
\title{2023 Homework 2\\Complete solutions and executable verification}
\author{EE3032 Introduction to VLSI Design}\date{13 October 2026}
\begin{document}\maketitle
The original two-page assignment is preserved in the assignments directory.
The three questions are analytical. The accompanying Python and SPICE files
provide additional reproducible verification without replacing the derivations.

\section*{1. Multi-threshold CMOS design}
The upper gate in Fig.~1 is XOR, preceded on its $A$ input by an inverter.
Consequently
\[
 U=\bar A,\quad V=\bar B,\qquad
 \boxed{D=U\oplus B=AB+\bar A\bar B},\qquad
 \boxed{E=\overline{BC}}.
\]
With the stated one-unit low-$V_t$ gate delay, the original paths are
$A\to\mathrm{INV}\to\mathrm{XOR}\to D$ (two units),
$B\to\mathrm{XOR}\to D$ (one unit), and
$B,C\to\mathrm{NAND}\to E$ (one unit). The critical delay is two units.
Use low-$V_t$ devices in the $A$ inverter and all devices realizing $D$.
Use high-$V_t$ devices in the four-transistor NAND for $E$, provided its
characterized delay is at most two units. This spends the NAND's one unit
of slack without extending the circuit's maximum delay.

The following complete static CMOS implementation generates both complements
locally. Its XOR core uses the pull-down condition $UB+AV$; the pull-up is
the exact series/parallel dual. All shown thin channels are low-$V_t$;
the four bold NAND channels are high-$V_t$. All pMOS bulks connect to
$\VDD$, and all nMOS bulks to ground. No body-terminal connections are implicit
internal-node ties.
\begin{center}
\resizebox{.98\linewidth}{!}{
\begin{tikzpicture}[font=\small]
 % Inverters for complements.
 \MOS{pa}{1}{4}{A}{p}{.5pt}\MOS{na}{1}{2.4}{A}{n}{.5pt}
 \draw (pat)--(1,4.9) node[above]{$\VDD$};\draw (pab)--(nat);
 \draw (1,3.2)--(2,3.2) node[right]{$U=\bar A$};
 \draw (nab)--(1,1.5) node[below]{$0$};
 \MOS{pb}{4}{4}{B}{p}{.5pt}\MOS{nb}{4}{2.4}{B}{n}{.5pt}
 \draw (pbt)--(4,4.9) node[above]{$\VDD$};\draw (pbb)--(nbt);
 \draw (4,3.2)--(5,3.2) node[right]{$V=\bar B$};
 \draw (nbb)--(4,1.5) node[below]{$0$};
 \node at (2.5,.6) {Complement generators: low-$V_t$};
 % XOR(U,B), equivalent to XNOR(A,B).
 \MOS{p0}{8}{5.5}{U}{p}{.5pt}\MOS{p1}{10.4}{5.5}{B}{p}{.5pt}
 \MOS{p2}{8}{4.2}{A}{p}{.5pt}\MOS{p3}{10.4}{4.2}{V}{p}{.5pt}
 \draw (p0t)--(p1t);\draw (9.2,5.95)--(9.2,6.5) node[above]{$\VDD$};
 \draw (p0b)--(p1b);\draw (p2t)--(p3t);\draw (9.2,5.05)--(9.2,4.65);
 \draw (p2b)--(p3b);\draw (9.2,3.75)--(9.2,3.2)--(11.4,3.2) node[right]{$D$};
 \MOS{n0}{8}{2.4}{U}{n}{.5pt}\MOS{n1}{8}{1.1}{B}{n}{.5pt}
 \MOS{n2}{10.4}{2.4}{A}{n}{.5pt}\MOS{n3}{10.4}{1.1}{V}{n}{.5pt}
 \draw (n0t)--(n2t);\draw (9.2,2.85)--(9.2,3.2);
 \draw (n0b)--(n1t);\draw (n2b)--(n3t);\draw (n1b)--(n3b);
 \draw (9.2,.65)--(9.2,.2) node[below]{$0$};
 \node at (9.2,-.7) {$D$: low-$V_t$ core};
 % NAND E.
 \MOS{ep0}{14}{5}{B}{p}{1.7pt}\MOS{ep1}{16.4}{5}{C}{p}{1.7pt}
 \draw (ep0t)--(ep1t);\draw (15.2,5.45)--(15.2,6) node[above]{$\VDD$};
 \draw (ep0b)--(ep1b);\draw (15.2,4.55)--(15.2,3.6)--(16.8,3.6) node[right]{$E$};
 \MOS{en0}{15.2}{2.8}{B}{n}{1.7pt}\MOS{en1}{15.2}{1.4}{C}{n}{1.7pt}
 \draw (15.2,3.6)--(en0t);\draw (en0b)--(en1t);
 \draw (en1b)--(15.2,.4) node[below]{$0$};
 \node at (15.2,-.5) {$E$: high-$V_t$ NAND};
\end{tikzpicture}}
\end{center}
The added $B$ complement generator is kept low-$V_t$: in this specific static
implementation the $B\to V\to D$ path also takes two gates. Merely selecting
all $B$-controlled XOR devices as high-$V_t$ would be unjustified because
those same series devices can carry a sensitized critical $A$ transition.
The prompt gives no high-$V_t$ delay or leakage parameters, so an unconditional
numerical timing guarantee or leakage reduction cannot be calculated.
If the high-$V_t$ NAND exceeds two units, retain low-$V_t$ devices or resize
and characterize it. The standard CIC018 TT model is not a dual-$V_t$ model.
\texttt{reproduce.py --logic-only} checks all eight $(A,B,C)$ assignments.

\section*{2. Non-ideal effect}
\subsection*{(a) Precharge}
The on PMOS raises BL to $\VDD$. The access NMOS gate is at $\VDD$;
charging the storage capacitance stops in the stated cutoff approximation
when $V_{GS}=\VDD-V_C=V_t$. Thus
\[
 \boxed{V_C(0)=\VDD-V_t.}
\]
The PMOS does not lose one NMOS threshold. If body effect is included,
$V_t$ is the effective threshold at the final source voltage, so the equation
is implicit rather than a subtraction of the zero-body-bias threshold.
\subsection*{(b) Retention}
With the switches off and the storage junction leaking a constant $I_L>0$,
\[
 C\frac{dV_C}{dt}=-I_L,\qquad
 V_C(t)=\VDD-V_t-\frac{I_L}{C}t.
\]
Set this equal to the buffer threshold $\VDD/2$:
\[
 \boxed{t_{\rm ret}=\frac{C(\VDD/2-V_t)}{I_L},\quad V_t<\VDD/2.}
\]
If $V_t\geq\VDD/2$, there is no initially valid stored high with positive
margin. For zero leakage and positive initial margin this ideal model has
infinite retention. The result includes the capacitance $C$, necessary
dimensionally; the Lin answer omits it. The Liao and Jian answers include it.

In the printed topology the buffer observes BL, on the other side of the
off access NMOS. Therefore this is the intended \emph{stored-charge}
retention, assuming an ideal nondestructive read. The continuously observed
Out while both switches remain off cannot be inferred from $C$ alone.
BL capacitance, leakage and the read protocol are missing. Upon reconnecting
an initially charged BL, even an ideal switch would give charge sharing
$V_{\rm share}=(CV_C+C_{BL}V_{BL,0})/(C+C_{BL})$.

For a clearly labelled example, take $\VDD=1.8$ V, $V_t=0.45$ V,
$C=20$ fF, $I_L=1$ pA: initial voltage $1.35$ V, switching boundary
$0.90$ V, and retention $9.00$ ms. These values are illustrative, not specified
assignment data. \texttt{q2\_retention.sp} implements the ideal capacitor
and current source; it is separate from the process transistor model.

\section*{3. Pseudo-NMOS NOR gate}
\subsection*{(a) Function}
The PMOS gate is grounded and its pull-up is always on. Two parallel NMOS
devices pull down when either $A$ or $B$ is high. Therefore
\[
 \boxed{V_o:\ \overline{A+B}.}
\]
This is ratioed logic with nonzero static current when the output is low.
For an arbitrary reference aspect ratio $S_0$, let
$s_p=(W/L)_p/S_0$, $s_A=(W/L)_{nA}/S_0$ and $s_B=(W/L)_{nB}/S_0$.
The stipulated resistance model gives
\[
 R_P=2R_n/s_p,\quad R_A=R_n/s_A,\quad R_B=R_n/s_B.
\]
\begin{center}\begin{tabular}{cccc}
\toprule $A$&$B$&Pull-down&$V_o/\VDD$\\\midrule
0&0&open&1\\
1&0&$R_A$&$s_p/(s_p+2s_A)$\\
0&1&$R_B$&$s_p/(s_p+2s_B)$\\
1&1&$R_A\parallel R_B$&$s_p/[s_p+2(s_A+s_B)]$\\\bottomrule
\end{tabular}\end{center}
\subsection*{(b) The wording changes the answer}
Unlike 2020 and 2021, the original 2023 PDF says \emph{smallest} logic-zero
voltage is $0.1\VDD$. This is not silently replaced by ``highest.''
For symmetric $r=s_A/s_p=s_B/s_p$, the smallest low is the two-on case:
\[
 V_{11}/\VDD=\frac1{1+4r}.
\]
If ``smallest is'' means equality, the literal answer is
\[
 \boxed{r=\frac94=2.25},\qquad
 V_{01}=V_{10}=\frac{2}{11}\VDD\simeq0.1818\VDD,
 \quad V_{11}=0.1\VDD.
\]
These remain below a $0.5\VDD$ decision boundary, but a receiver threshold
is not specified in Question~3. If ``smallest'' means a lower bound
$V_{11}\geq0.1\VDD$, it gives $r\leq9/4$ and does not uniquely determine size.
For unequal NMOS sizes the equality gives $s_A+s_B=(9/2)s_p$.

If the intended requirement instead matches the earlier years' conventional
\emph{maximum acceptable} low $V_{OL,\max}=0.1\VDD$, each one-on case must
satisfy $s_p/(s_p+2s_A)\leq0.1$ and its $B$ counterpart. Then
\[
 \boxed{(W/L)_{nA},(W/L)_{nB}\geq\frac92(W/L)_p.}
\]
At the symmetric boundary $r=4.5$, $V_{01}=V_{10}=0.1\VDD$ and
$V_{11}=\VDD/19$. These are ratios; absolute widths and lengths are free.
Lin and Liao use this maximum-low interpretation. Jian discusses a lower
bound for the two-on state, but its inequalities do not identify the same
conventional sizing condition. Both interpretations are explicitly separated
here instead of merging incompatible conditions.

\clearpage
\section*{Executable checks and reference comparison}
\texttt{sizing.csv} records the two literal/maximum boundary ratios and an
additional margin example. \texttt{q3\_resistance.sp} measures all four
divider cases at the conventional boundary ($R_P=9$ k$\Omega$, each NMOS
resistance $1$ k$\Omega$). Its predicted voltages at $1.8$ V are
$(1.8,0.18,0.18,1.8/19)$ V. \texttt{q3\_cic018.sp} separately sweeps
the actual BSIM3 devices at ratios 1, 2, 2.25, 3, 4.5, 4, 6 and 8.
The order of the sweep is immaterial. BSIM3 does not enforce constant
$R_P=2R_n$; its voltages must not substitute for the stipulated divider answer.

The three supplied 2023 HW2 submissions (Lin 110501558, Liao 110501512,
Jian 109501570) were read before derivation and compared above. Their
threshold-retention model and NOR truth table are retained where valid.
The complete transistor drawing, explicit dual-$V_t$ timing condition,
correct dimensional retention formula, BL qualification and the separate
literal sizing interpretation improve the combined answer.

Run \texttt{python reproduce.py --work PATH} to regenerate the logic checks,
three real HSPICE measurement tables and \texttt{validation.json}.
Use \texttt{--model PATH} and \texttt{--hspice PATH} for other installations.
The process file remains local. All stored measurements are actual HSPICE
outputs; analytical checks are labelled separately.

\subsection*{Recorded execution}
All three decks completed with PrimeSim HSPICE S-2021.09, with exit code zero.
The eight logic cases passed. Ideal retention measured $9.000025$ ms versus
$9.000000$ ms analytically; the $25$ ns difference is startup integration,
about $2.78\times10^{-6}$ relative error. The four resistor results match
the predicted values to the saved precision. The separate TT BSIM3 sweep
uses $W_p=0.5\,\mu$m, $L_n=L_p=0.18\,\mu$m and $W_n=rW_p$:
\begin{center}\begin{tabular}{rrrr}
\toprule $r$&$V_{01}=V_{10}$ (V)&$V_{11}$ (V)&$V_{00}$ (V)\\\midrule
1&0.176482&0.078389&1.800000\\
2&0.078122&0.036908&1.800000\\
2.25&0.068486&0.032568&1.800000\\
3&0.049900&0.024040&1.800000\\
4.5&0.032241&0.015732&1.800000\\
4&0.036566&0.017786&1.800000\\
6&0.023779&0.011676&1.800000\\
8&0.017604&0.008684&1.800000\\\bottomrule
\end{tabular}\end{center}
These voltages quantify the chosen process-model example, not the exact
constant-resistance size boundary. The transistor deck also saves supply
currents and static power. The raw listings remain local; the compact CSV
tables and reproduction scripts are included in the repository.
\end{document}
